% !PRINTED 236
\weberpage{236}
%<<CONT
\[
\begin{aligned}
5\,\sqrt{c}^{\,4} - 4\,c_1\,\sqrt{c}^{\,3} + 3\,c_2\,\sqrt{c}^{\,2} - 2\,c_3\,\sqrt{c} + c_4 &= 0,\\
-10\,\sqrt{c}^{\,3} + 6\,c_1\,\sqrt{c}^{\,2} - 3\,c_2\,\sqrt{c} + c_3 &= 0,
\end{aligned}
\]
and on account of the double sign of $\sqrt{c}$ these equations split into the four
\[
\begin{aligned}
5\,c^2 + 3\,c_2\,c + c_4 &= 0,\\
4\,c_1\,c + 2\,c_3 &= 0,\\
10\,c + 3\,c_2 &= 0,\\
6\,c_1\,c + c_3 &= 0.
\end{aligned}
\]
From the second and fourth of these equations it follows that
\[
c_1 = 0, \quad c_3 = 0,
\]
and then from the third and first
\[
c_2 = -\frac{10}{3}\,c, \quad c_4 = -5\,c^2,
\]
so that the equation for $y$ thus takes the form
\weq{10}{y^5 - \frac{10}{3}\,c\,y^3 + 5\,c^2\,y + c_5 = 0.}
This equation still depends on the two parameters $c_5$ and $c$; one can reduce it to an equation with one parameter by the substitution
\weq{11}{y = \sqrt{\frac{c}{3}}\,z, \quad \gamma = c_5\,\sqrt{\frac{3^5}{c^5}},}
whereby it takes the simple and elegant form
\weq{12}{z^5 - 10\,z^3 + 45\,z + \gamma = 0.}
This is the \gesperrt{Brioschi} normal form. The substitution (11), however, suffers from the drawback that it still contains a square root; whereas in the formulas (8) and (10) only the two square roots discussed in the previous section occur. But, following a remark of \gesperrt{Klein}, one can also pass from (10) by rational means to a normal form that contains only one parameter.

If, for example, one sets
\weq{13}{y = \frac{3\,c_5}{c^2}\,\frac{1}{z}, \quad \gamma = \frac{9\,c_5^2}{c^5},}
then from (10) one obtains the equation
\weq{14}{z^5 + 15\,z^4 - 10\,\gamma\,z^2 + 3\,\gamma^2 = 0.}
